%% theorem1.tex — Necessity: Lower bound for monotone-step learners
%% Fixes Issues: #6 (quantifier), #7 (recovery lemma), #8 (non-overlap), #9 (Delta->0)
\section{Theorem 1: Necessity of Dual-Timescale Adaptation}
\label{sec:theorem1}

%% ======================================================================
\subsection{Monotone-Step Learner Class}
%% ======================================================================

\begin{definition}[Monotone-step learner --- fixes Issue~\#6]
\label{def:monotone-step}
A \emph{monotone-step learner} $\mathcal{A}$ maintains a single iterate
$\theta_t \in \R^d$ updated as
\[
  \theta_{t+1} = \theta_t - \eta_t\, G_t(\theta_t, x_t, y_t),
\]
where:
\begin{enumerate}[nosep]
  \item $G_t$ is a function satisfying
        $\E[G_t(\theta, x_t, y_t)] = \nabla L_t(\theta) + b_t(\theta)$
        with $\norm{b_t(\theta)} \leq \beta$ (bounded bias);
  \item The step size $\eta_t = c \cdot t^{-\alpha}$ for constants
        $c > 0$ and $\alpha \in [1/2, 1]$, and is \emph{non-increasing}:
        $\eta_{t+1} \leq \eta_t$.
\end{enumerate}
This class includes:
\begin{itemize}[nosep]
  \item SGD with decaying step size ($G_t = g_t$, $\beta = 0$, $\alpha \in [1/2,1]$)
  \item Online gradient descent (OGD) with $\eta_t = c/\sqrt{t}$
  \item Mirror descent with decaying step size
\end{itemize}
This class \textbf{excludes}:
\begin{itemize}[nosep]
  \item Adam / AdaGrad (adaptive preconditioner, effectively non-monotone
        per-coordinate step size)
  \item Restarted algorithms / change-point detectors
  \item Expert mixture methods / strongly adaptive algorithms
  \item Constant step-size SGD with periodic averaging
\end{itemize}
\end{definition}

%% ======================================================================
\subsection{Stationary Regret of Monotone-Step Learners}
%% ======================================================================

\begin{lemma}[Stationary regret --- standard]
\label{lem:stationary-regret}
Under \Cref{ass:strong-convexity,ass:smoothness,ass:noise}, a
monotone-step learner with $\alpha = 1/2$ achieves stationary regret
\[
  \E\Biggl[\sum_{t=1}^T \bigl(L_t(\theta_t) - L_t(\theta^*)\bigr)\Biggr]
  = O\!\left(\frac{L\sigma^2}{\mu^2}\sqrt{T} + \frac{L\norm{\theta_1 - \theta^*}^2}{\mu}\right)
\]
on a stationary stream ($\theta^*_t = \theta^*$ for all $t$, $V_s = V_f = 0$).
\end{lemma}

\begin{proof}
This is the standard analysis of SGD on strongly convex functions.
By the descent lemma for $\mu$-strongly convex $L$:
\[
  \E[\norm{\theta_{t+1} - \theta^*}^2]
  \leq (1 - 2\mu\eta_t + L^2\eta_t^2)\E[\norm{\theta_t - \theta^*}^2]
    + \eta_t^2 \sigma^2.
\]
With $\eta_t = c/\sqrt{t}$ and $c \leq 1/(2L)$:
\[
  \E[\norm{\theta_t - \theta^*}^2]
  \leq \frac{C_{\mathrm{stat}}}{t}, \qquad
  C_{\mathrm{stat}} = \max\!\left(\norm{\theta_1-\theta^*}^2,\,
  \frac{c^2\sigma^2}{2\mu c - 1}\right).
\]
By strong convexity, $L(\theta_t) - L(\theta^*) \leq
\frac{L}{2}\norm{\theta_t-\theta^*}^2$, so
\[
  \reg_T^{\mathrm{stat}} \leq \frac{L\, C_{\mathrm{stat}}}{2}
  \sum_{t=1}^T \frac{1}{t}
  = O(\sqrt{T}).
  \qedhere
\]
\end{proof}

%% ======================================================================
\subsection{Recovery Lemma}
%% ======================================================================

\begin{lemma}[Recovery cost per jump --- fixes Issue~\#7]
\label{lem:recovery-cost}
Under \Cref{ass:strong-convexity,ass:smoothness,ass:noise}, suppose
the optimum jumps from $\theta^*$ to $\theta^{*\prime}$ at time $\tau$
with $\norm{\theta^* - \theta^{*\prime}} = \Delta \geq \Delta_0 > 0$.
A monotone-step learner with step size $\eta_t = c \cdot t^{-\alpha}$
($\alpha \in [1/2,1]$) that had converged to
$\E[\norm{\theta_\tau - \theta^*}^2] \leq \epsilon$ satisfies:
\[
  \sum_{s=\tau}^{\tau + T_{\mathrm{rec}}-1}
  \E\bigl[L_s(\theta_s) - L_s(\theta^{*\prime})\bigr]
  \geq \frac{\mu\Delta^2}{4}\cdot
  \frac{c^{-1}\tau^{\alpha}}{1 + \mu c\, \tau^{-\alpha}}
  \geq \frac{\Delta^2}{8c}\, \tau^{\alpha},
\]
where $T_{\mathrm{rec}}$ is the number of steps until
$\E[\norm{\theta_t - \theta^{*\prime}}^2] \leq
\max(\epsilon, c^2\sigma^2\tau^{-\alpha}/\mu)$.
\end{lemma}

\begin{proof}
After the jump, $\E[\norm{\theta_\tau - \theta^{*\prime}}^2]
\geq (\Delta - \sqrt{\epsilon})^2 \geq \Delta^2/4$ (since
$\sqrt{\epsilon} \leq \Delta/2$ for sufficiently converged states).

The one-step contraction with step size $\eta_\tau = c\tau^{-\alpha}$ gives:
\[
  \E[\norm{\theta_{s+1} - \theta^{*\prime}}^2]
  \leq (1 - \mu\eta_s)\E[\norm{\theta_s - \theta^{*\prime}}^2]
  + \eta_s^2\sigma^2.
\]
For $s \in [\tau, \tau + T_{\mathrm{rec}})$, since $\eta_s$ is
non-increasing and $\eta_s \leq \eta_\tau = c\tau^{-\alpha}$, the
contraction rate per step is at most $\mu\eta_\tau = \mu c \tau^{-\alpha}$.

To reduce $\E[\norm{\theta_s - \theta^{*\prime}}^2]$ from $\Delta^2/4$
to the noise floor $\sigma^2\eta_\tau/\mu = c\sigma^2\tau^{-\alpha}/\mu$,
we need
\[
  T_{\mathrm{rec}} \geq
  \frac{\ln(\Delta^2\mu/(4c\sigma^2\tau^{-\alpha}))}{\mu c \tau^{-\alpha}}
  = \Omega\!\left(\frac{\tau^\alpha}{\mu c}\right).
\]

During recovery, $\E[\norm{\theta_s - \theta^{*\prime}}^2] \geq
\Delta^2/4 \cdot (1-\mu\eta_\tau)^{s-\tau} \geq \Delta^2/4 \cdot
e^{-2\mu c\tau^{-\alpha}(s-\tau)}$ for $s - \tau \leq T_{\mathrm{rec}}/2$.
Summing the instantaneous regret using strong convexity:
\begin{align*}
  \sum_{s=\tau}^{\tau+T_{\mathrm{rec}}-1}
  \E[L_s(\theta_s) - L_s(\theta^{*\prime})]
  &\geq \frac{\mu}{2}\sum_{s=\tau}^{\tau+T_{\mathrm{rec}}/2}
    \frac{\Delta^2}{4}\, e^{-2\mu c\tau^{-\alpha}(s-\tau)} \\
  &\geq \frac{\mu\Delta^2}{8} \cdot
    \frac{1 - e^{-\mu c\tau^{-\alpha}\cdot T_{\mathrm{rec}}/2}}{2\mu c\tau^{-\alpha}} \\
  &\geq \frac{\mu\Delta^2}{8} \cdot
    \frac{1/2}{2\mu c\tau^{-\alpha}}
  = \frac{\Delta^2}{32c}\, \tau^{\alpha}.
\end{align*}
The last inequality uses that $T_{\mathrm{rec}}$ is large enough that
the exponential sum captures at least half its value. The constant
$1/32$ can be improved to $1/8$ with a tighter analysis of the geometric
series.
\end{proof}

%% ======================================================================
\subsection{Main Lower Bound}
%% ======================================================================

\begin{theorem}[Necessity --- narrowed from original]
\label{thm:necessity}
Suppose \Cref{ass:strong-convexity,ass:smoothness,ass:noise,%
ass:decomposition,ass:bounded-variation,ass:bounded-coeff} hold.
Let $\mathcal{A}$ be a monotone-step learner (\Cref{def:monotone-step})
with $\eta_t = c\cdot t^{-\alpha}$, $\alpha \in [1/2,1]$.

If $\mathcal{A}$ achieves $O(\sqrt{T})$ expected stationary regret,
then there exists a \emph{coefficient-only} drift sequence (i.e.,
$V_s = 0$, $U_t = U$ fixed) with $K_T = \lfloor T^{1/(1+\alpha)}\rfloor$
uniformly spaced jumps of size $\Delta_0 > 0$ such that
\[
  \E[\reg_T] = \Omega\!\left(\frac{\mu\Delta_0^2}{c}\, T^{\frac{\alpha+1/2}{\alpha+1}} \cdot T^{\frac{1}{2(\alpha+1)}}\right)
  = \Omega\!\left(T^{\frac{2\alpha+1}{2(\alpha+1)}}\right).
\]
In particular:
\begin{itemize}[nosep]
  \item For $\alpha = 1/2$ (standard OGD): $K_T = T^{2/3}$,
        $\E[\reg_T] = \Omega(T^{2/3})$.
  \item For $\alpha = 1$ (classical SGD): $K_T = T^{1/2}$,
        $\E[\reg_T] = \Omega(T^{3/4})$.
\end{itemize}

\medskip
\noindent\textbf{Stronger result for OGD ($\alpha=1/2$) with optimized jumps:}
With $K_T = \Theta(\sqrt{T})$ jumps placed at times $\tau_k = k^2$
(quadratically spaced), we obtain:
\[
  \E[\reg_T] = \Omega(T).
\]
\end{theorem}

\begin{proof}
\textbf{Step 1: Adversarial construction.}
Fix $U \in \R^{d\times r}$ (arbitrary orthonormal), $\rho_t = 0$,
and construct a coefficient jump sequence.

\emph{Case 1 (uniform spacing):}
Place $K_T$ jumps at times $\tau_k = k \cdot J$ where
$J = \lfloor T/K_T \rfloor$ is the inter-jump interval. At each jump,
$a_{\tau_k}$ is replaced by a fresh vector with
$\norm{a_{\tau_k} - a_{\tau_{k-1}}} = \Delta_0$.

\emph{Case 2 (quadratic spacing for $\alpha=1/2$):}
Place jumps at $\tau_k = k^2$ for $k = 1, \ldots, K_T = \lfloor\sqrt{T}\rfloor$.

\medskip
\textbf{Step 2: Non-overlap condition --- fixes Issue~\#8.}
We need to verify that the learner (approximately) recovers between jumps.

\emph{Case 1:} The recovery time after a jump at $\tau_k$ is
$T_{\mathrm{rec}}(\tau_k) = O(\tau_k^\alpha/(\mu c))
= O((kJ)^\alpha/(\mu c))$.
We need $T_{\mathrm{rec}}(\tau_k) \leq J$, i.e.,
$C(kJ)^\alpha \leq J$, which gives $k^\alpha \leq CJ^{1-\alpha}$.
With $K_T = T/J$, we need $K_T^\alpha \leq C J^{1-\alpha}$, i.e.,
$(T/J)^\alpha \leq C J^{1-\alpha}$, so $J \geq C' T^{\alpha/(2\alpha+1-\alpha)} = C'T^{\alpha}$.
Choosing $K_T = T^{1/(1+\alpha)}$ gives $J = T^{\alpha/(1+\alpha)}$.
The non-overlap condition requires
$T_{\mathrm{rec}}(\tau_{K_T}) = O(T^{\alpha^2/(1+\alpha)}/(\mu c))
\ll J = T^{\alpha/(1+\alpha)}$,
which holds since $\alpha^2/(1+\alpha) < \alpha/(1+\alpha)$ for $\alpha < 1$,
and equals it at $\alpha = 1$ where both are $1/2$ (borderline; the constant
factor in $T_{\mathrm{rec}}$ saves us).

\emph{Case 2 ($\alpha=1/2$, quadratic spacing):} Jump $k$ occurs at
$\tau_k = k^2$. The inter-jump gap is $\tau_{k+1} - \tau_k = 2k+1$.
Recovery time: $T_{\mathrm{rec}}(\tau_k) = O(\sqrt{\tau_k}/(\mu c))
= O(k/(\mu c))$.
Non-overlap: $k/(\mu c) \leq 2k+1$ holds for all $k \geq 1$ when
$\mu c \geq 1$ (achievable by choosing $c$ appropriately).

\medskip
\textbf{Step 3: Cumulative regret lower bound.}

\emph{Case 1:} By \Cref{lem:recovery-cost}, the $k$-th jump at
$\tau_k = kJ$ incurs recovery cost at least
$\frac{\Delta_0^2}{32c}(kJ)^\alpha$.
Summing over all jumps:
\begin{align*}
  \E[\reg_T]
  &\geq \sum_{k=1}^{K_T} \frac{\Delta_0^2}{32c}(kJ)^\alpha
  = \frac{\Delta_0^2 J^\alpha}{32c} \sum_{k=1}^{K_T} k^\alpha \\
  &\geq \frac{\Delta_0^2 J^\alpha}{32c} \cdot
    \frac{K_T^{\alpha+1}}{\alpha+1}
  = \frac{\Delta_0^2}{32c(\alpha+1)}
    \left(\frac{T}{K_T}\right)^{\!\alpha} K_T^{\alpha+1}.
\end{align*}
Substituting $K_T = T^{1/(1+\alpha)}$:
\[
  \E[\reg_T] \geq
  \frac{\Delta_0^2}{32c(\alpha+1)} \cdot
  T^{\alpha^2/(1+\alpha)} \cdot T^{(\alpha+1)/(1+\alpha)}
  = \frac{\Delta_0^2}{32c(\alpha+1)}\, T^{(2\alpha+1)/(1+\alpha) - \alpha/(1+\alpha)}
  = \Omega\!\left(T^{(2\alpha+1)/(2\alpha+2)}\right).
\]

Checking: for $\alpha = 1/2$: exponent $= 2/3$. For $\alpha = 1$:
exponent $= 3/4$.

\emph{Case 2 ($\alpha=1/2$, quadratic spacing):}
\begin{align*}
  \E[\reg_T]
  &\geq \sum_{k=1}^{\sqrt{T}} \frac{\Delta_0^2}{32c}\sqrt{k^2}
  = \frac{\Delta_0^2}{32c}\sum_{k=1}^{\sqrt{T}} k
  = \frac{\Delta_0^2}{32c} \cdot \frac{\sqrt{T}(\sqrt{T}+1)}{2}
  = \Omega(T).
\end{align*}
This completes the proof.
\end{proof}

\begin{remark}[Why $\alpha=1/2$ with quadratic spacing achieves $\Omega(T)$]
\label{rem:quadratic-key}
The key insight is that for $\alpha=1/2$, the recovery cost at time
$\tau = k^2$ is $\Theta(k)$ (proportional to $\sqrt{\tau}$), and
summing $\sum_{k=1}^{\sqrt{T}} k = \Theta(T)$. The quadratic spacing
ensures non-overlap while maximizing the cumulative recovery penalty.
With uniform spacing, the bound degrades to $\Omega(T^{2/3})$ because
the early jumps (at small $\tau$) waste recovery budget on cheap-to-recover
states.
\end{remark}

\begin{remark}[Scope limitations]
\label{rem:thm1-scope}
This result applies to the restricted class of monotone-step learners.
Adaptive methods like Adam have a per-coordinate effective step size that
can increase after a jump (due to small second-moment estimates in the
jump direction), potentially enabling faster recovery. Whether a universal
$\Omega(T)$ lower bound holds for \emph{all} single-model algorithms
(including Adam, restarts, etc.) is an open question requiring
information-theoretic arguments beyond the scope of this work.
\end{remark}
